\begin{foldable} % Overview
  We evaluate DIP-KD on extensive BBDFKD benchmarks, supplemented by ablation studies to validate its mechanics and practical effectiveness.
\end{foldable}

\subsection{Experimental setup} \label{ssec:04-experiments-setup}
\begin{foldable} % Datasets, architectures, settings
  We evaluate our method across 12 benchmarks, including 8 general-purpose datasets (USPS, MNIST, SVHN, FMNIST, CIFAR10, CIFAR100, Tiny-ImageNet, Imagenette)~\cite{lecun2002gradient,hull1994database,netzer2011reading,xiao2017fashion,krizhevsky2009learning,le2015tiny,howard2020imagenette} and 4 subsets from the medical dataset MedMNIST~\cite{yang2021medmnist}.
  We investigate three teacher-student architectures: LeNet5~\cite{lecun2002gradient}, AlexNet~\cite{krizhevsky2012imagenet}, and ResNet~\cite{he2016deep}.
  These datasets and architectures are the de facto evaluation in BBDFKD methods~\cite{wang2021zero, zhang2022ideal, yuan2024data}.
\end{foldable}

\subsection{Baselines} \label{ssec:04-experiments-baselines}
\begin{foldable} % Baselines
  We compare our method DIP-KD with the baselines:
  \begin{itemize}
    \item \textbf{NaiveKD}: The student is trained on uniform noise $x \sim \mathcal{X}_\mathcal{U}$, where $\mathcal{X}_\mathcal{U}=\mathcal{U}[0,1]^{C \times H \times W}$ and the teacher's hard labels $\hat{y}_{T}^\text{hard} = T(x) \in \{1, \ldots, K\}$.
    \item \textbf{BBDFKD methods}: We reproduce three state-of-the-art methods: ZSDB3~\cite{wang2021zero}, IDEAL~\cite{zhang2022ideal}, and DFHL-RS~\cite{yuan2024data}.
  \end{itemize}
  For fair comparison, we use the same teacher-student network architectures and synthetic image budget for all methods.
  We report the \textit{mean} $\pm$ \textit{standard error} (SE) of distillation accuracy, which is the student accuracy on the hold-out test set, across five runs.
  The accuracies of baseline BBDFKD methods are reproduced with their official implementations\footnote{ZSDB3 at \url{https://github.com/zwang84/zsdb3kd}}\,\footnote{IDEAL at \url{https://github.com/SonyResearch/IDEAL}}\,\footnote{DFHL-RS at \url{https://github.com/LetheSec/DFHL-RS-Attack}}.
  We ensured that all baseline implementations were tuned ethically and fairly following the recommended settings.
\end{foldable}

\subsection{Standard experiments} \label{ssec:04-experiments-standard}
\begin{foldable} % Overview
  We categorize the benchmarks as \textit{simple}, \textit{complex}, and \textit{domain-specific} datasets.
  Four simple datasets---MNIST, USPS, SVHN, FMNIST---all have 10 classes of simple objects, at a $32\times32$ or smaller resolution.
  The complex datasets have more complicated objects or more classes: CIFAR10 and CIFAR100 (10 and 100 classes, $32\times32$), Tiny-ImageNet (200 classes, $64\times64$), and Imagenette (10 classes, but high-resolution).
  Finally, we evaluate the domain-specific performance on 4 subsets of MedMNIST, including BreastMNIST, OrganAMNIST, DermaMNIST, and BloodMNIST, which include $32\times32$ images of varying medical modalities.
  We report the distillation accuracy of all methods in Table~\ref{tab:standard-all}, where our method \textbf{DIP-KD} consistently outperforms all baselines across 11/12 datasets, with significant margins on complex and domain-specific ones.
\end{foldable}

\begin{table}[ht]
  \centering
  \begin{threeparttable}
    \caption{Distillation accuracy (\%) on simple, complex, and domain-specific datasets.}
    \label{tab:standard-all}
    \small
    \begin{tabular}{l c c c c c c c}
      \toprule
      \boldtight{Dataset} & \boldtight{Setup} & \boldtight{Teacher} & \boldtight{NaiveKD} & \boldtight{ZSDB3} & \boldtight{IDEAL} & \boldtight{DFHL-RS} & \boldtight{DIP-KD} \\
      \midrule

      \multicolumn{8}{c}{\textbf{Simple datasets}} \\
      \midrule
      MNIST    & L, 20K  & 99.28 & 96.54$_{\pm 0.60}$           & 96.66$_{\pm 0.39}$           & 96.96$_{\pm 0.34}$ & 98.53$_{\pm 0.10}$          & \textbf{98.59}$_{\pm 0.08}$ \\
      USPS     & L, 50K  & 95.47 & 42.05$_{\pm 0.08}$           & 65.50$_{\pm 0.16}$           & 92.66$_{\pm 0.57}$ & 93.97$_{\pm 0.60}$          & \textbf{94.20}$_{\pm 0.21}$ \\
      SVHN     & A, 50K  & 96.16 & 83.35$_{\pm 0.57}$           & 85.38$_{\pm 1.69}$           & 87.86$_{\pm 0.93}$ & 95.41$_{\pm 0.13}$          & \textbf{95.94}$_{\pm 0.05}$ \\
      FMNIST   & L, 50K  & 90.90 & 55.11$_{\pm 1.23}$           & 71.60$_{\pm 2.35}$           & 82.84$_{\pm 1.22}$ & 83.89$_{\pm 1.35}$          & \textbf{83.98}$_{\pm 0.65}$ \\
      \midrule

      \multicolumn{8}{c}{\textbf{Complex datasets}} \\
      \midrule
      CIFAR10  & R, 50K  & 95.21 & 13.99$_{\pm 0.23}$           & 56.74$_{\pm 1.98}$           & 67.58$_{\pm 1.29}$ & 76.50$_{\pm 0.84}$          & \textbf{83.41}$_{\pm 0.46}$ \\
      CIFAR100 & R, 100K & 78.36 & \phantom{0}2.68$_{\pm 0.13}$ & \phantom{0}9.19$_{\pm 0.39}$ & 33.64$_{\pm 0.91}$ & 51.97$_{\pm 0.11}$          & \textbf{52.85}$_{\pm 0.30}$ \\
      TinyImgN & R, 500K & 66.08 & \phantom{0}0.75$_{\pm 0.07}$ & \phantom{0}3.29$_{\pm 0.16}$ & 23.76$_{\pm 1.03}$ & \textbf{30.52}$_{\pm 1.43}$ & 28.03$_{\pm 0.22}$ \\
      Imagenette & R, 50K  & 91.29 & 33.21$_{\pm 0.39}$         & 33.43$_{\pm 0.85}$           & 36.90$_{\pm 0.38}$ & 40.53$_{\pm 2.36}$          & \textbf{61.84}$_{\pm 0.77}$ \\
      \midrule

      \multicolumn{8}{c}{\textbf{Domain-specific datasets}} \\
      \midrule
      BreastM  & L, 50K  & 79.49 & 61.78$_{\pm 1.60}$           & 63.37$_{\pm 1.48}$           & 67.54$_{\pm 1.27}$ & 73.81$_{\pm 1.80}$          & \textbf{77.95}$_{\pm 0.77}$ \\
      OrganAM  & L, 50K  & 87.60 & 24.09$_{\pm 0.44}$           & 51.81$_{\pm 1.21}$           & 69.09$_{\pm 1.92}$ & 75.12$_{\pm 2.89}$          & \textbf{82.32}$_{\pm 1.08}$ \\
      DermaM   & R, 50K  & 79.60 & 56.82$_{\pm 1.50}$           & 60.15$_{\pm 2.93}$           & 62.76$_{\pm 2.76}$ & 64.71$_{\pm 3.02}$          & \textbf{71.29}$_{\pm 2.67}$ \\
      BloodM   & R, 50K  & 98.19 & \phantom{0}9.09$_{\pm 0.13}$ & 39.62$_{\pm 1.67}$           & 55.28$_{\pm 2.18}$ & 61.67$_{\pm 3.11}$          & \textbf{73.98}$_{\pm 1.77}$ \\
      \bottomrule
    \end{tabular}
    \begin{tablenotes}[flushleft] \footnotesize
    \item {Setup denotes architecture \{\textbf{\underline{L}}eNet5, \textbf{\underline{A}}lexNet, \textbf{\underline{R}}esNet18\} and synthetic dataset size; [\textbf{bold}] best results.}
    \end{tablenotes}
  \end{threeparttable}
\end{table}

\begin{foldable}
  While performance saturates across all baselines on simple datasets, the accuracy gap widens as domain complexity scales.
  NaiveKD is first to collapse to near random-guessing levels in complex datasets, implying that unstructured noise cannot elicit the deep semantic knowledge embedded in black-box experts.
  Notably, on high-resolution images like Imagenette, our DIP-KD achieves a +21.31\% absolute gain over the strongest baseline, DFHL-RS.
  On domain-specific medical tasks, DIP-KD outperforms baselines by at least +4\% to +12\%.
  We also notice that abundant classes represent an interesting boundary condition, where DIP-KD remains highly competitive on CIFAR100 and approaches state-of-the-art on Tiny-ImageNet.
  These results suggest that data diversity is a critical factor for effective knowledge acquisition, particularly as task complexity and domain specificity increase.
\end{foldable}

\begin{foldable}
  In terms of robustness, we observe that DIP-KD has consistently low standard errors on most datasets, indicating strong stability across runs.
  We attribute this robustness to our integrated focus on diversity and signal restoration, which provide the student with softer training signals than noise-heavy or hard-label baselines.
\end{foldable}

\subsection{Ablation studies} \label{ssec:04-experiments-ablation}
\begin{foldable}
  In this section, we provide deeper insights into our framework through component dissection and practical application.
  We first investigate the contribution of individual components, the effect of synthetic data budget, and the semantics of our image priors.
  Then, we further extend our evaluation to real-world scenarios, including cross-architecture settings and aggressive model compression.
\end{foldable}

\subsubsection{Contribution of components} \label{sssec:04-experiments-ablation-components}

\begin{table}[ht]
  \centering
  \begin{threeparttable}
    \caption{Distillation accuracy (\%) by inclusion of components.}
    \label{tab:abla-components}
    \begin{tabular}{l l l}
      \toprule
      \makecell{\textbf{Method}} & \makecell{\textbf{MNIST}} & \makecell{\textbf{CIFAR10}} \\
      \midrule
      NaiveKD (noise $x \sim \mathcal{X}_\mathcal{U}$, Hard-KD) & \phantom{$+$ }\meanstd{96.02}{.60} & \phantom{$+$ }\meanstd{13.99}{.23} \\
      $+$ \textbf{Synthesis:}                                   & $+$ \phantom{0}2.27                & $+$ 64.18 \\
      \phantom{12} \bltrg Hierarchical noise                    & \phantom{$+$ 0}1.79\bltrg          & \phantom{$+$} 59.26\bltrg  \\
      \phantom{12} \bltrg Nonlinear transform                   & \phantom{$+$ 0}0.18\bltrg          & \phantom{$+$ 0}2.38\bltrg  \\
      \phantom{12} \bltrg Semantic cutmixing                    & \phantom{$+$ 0}0.30\bltrg          & \phantom{$+$ 0}2.54\bltrg  \\
      $+$ \textbf{Contrast:} Update $\mathcal{D}$ via $\mathcal{L}_\text{CT}$ & $+$ \phantom{0}0.16  & $+$ \phantom{0}3.14 \\
      $+$ \textbf{Distillation:} Add Soft-KD to $\mathcal{L}_S$ & $+$ \phantom{0}0.14                & $+$ \phantom{0}2.10 \\
      \midrule
      Complete \textbf{DIP-KD}                                  & \phantom{$+$ }98.59$_{\pm .08}$    & \phantom{$+$ }83.41$_{\pm .46}$ \\
      \bottomrule
    \end{tabular}
    \begin{tablenotes}[flushleft] \footnotesize
    \item {[$\blacktriangle$] denotes individual contribution of each Synthesis sub-component.}
    \end{tablenotes}
  \end{threeparttable}
\end{table}

\begin{foldable}
  We evaluate each component on MNIST and CIFAR10 to represent simple and complex benchmarks.
  Starting from NaiveKD, we sequentially add components and report student accuracy in Table~\ref{tab:abla-components}.
  Synthesis provides the largest gain ($+$2.27\% on MNIST, $+$64.18\% on CIFAR10), primarily driven by hierarchical noise.
  This is equivalent to the primer student $S_0$ trained on $\mathcal{X}_\mathcal{P}$, achieving 98\% and 78\% on MNIST and CIFAR10 (which already outperforms the baselines from Sec. \ref{ssec:04-experiments-standard}).
  This confirms that $S_0$ is sufficiently competent to perform contrastive optimization and soft distillation.
  For instance, on CIFAR10, we achieve $+$3.14\% via Contrast and $+$2.10\% via Distillation.
  These results confirm that the Synthesis pipeline plays a vital role in generating diverse data, while the novel use of the primer student for Contrast and Distillation phases significantly boosts KD performance beyond prior works.
\end{foldable}

\subsubsection{Embeddings analysis --- Synthesis} \label{sssec:04-experiments-ablation-embeddings-generation}
\begin{figure}[ht] % fig:embeddings-universal
  \centering
  \includegraphics[trim={120 0 0 0},clip,width=0.90\columnwidth]{./figures/04-experiments-vizembed-2d-objects.png}
  \caption{Feature embeddings of complex datasets, noise ($\mathcal{X}_\mathcal{U}$, $\mathcal{X}_\mathcal{N}$), and image priors $\mathcal{X}_\mathcal{P}$.}
  \label{fig:embeddings-universal}
\end{figure}

\begin{foldable}
  To assess the diversity and semantics of our image priors, we sample synthetic images and extract embeddings from a ResNet18 teacher backbone pre-trained on CIFAR10.\footnote{This is the 95.21\% ResNet18 teacher on CIFAR10, being reused as a held-out ``oracle''.}
  These 512-dimensional features are reduced to two dimensions using UMAP~\cite{mcinnes2018umap} with default configuration.
  We explore the remaining complex datasets CIFAR100, Tiny-ImageNet, and Imagenette, as well as uniform noise $\mathcal{X}_\mathcal{U}$, Gaussian noise $\mathcal{X}_\mathcal{N}$, and our image priors $\mathcal{X}_\mathcal{P}$.
  We declare that the use of real data and the teacher's backbone here is solely for evaluation purposes.
  We visualize these embeddings in Fig.~\ref{fig:embeddings-universal}.
\end{foldable}

\begin{table}[ht]
  \centering
  \begin{threeparttable}
    \caption{Distribution metrics (\%) on complex datasets.}
    \label{tab:abla-embeddings-coverage-objects}
    \small
    \begin{tabular}{c c c c c c}
      \toprule
      \textbf{Metric} & \textbf{Source} & \textbf{CIFAR10} & \textbf{CIFAR100} & \textbf{Tiny-ImageNet} & \textbf{Imagenette} \\
      \midrule
      & $\mathcal{X}_\mathcal{U}$ & \phantom{0}20.00 &           100.65 &           108.65 &           101.20 \\
      \textbf{Density}   & $\mathcal{X}_\mathcal{N}$ & \phantom{0}20.00 &           107.98 & \phantom{0}90.87 &           103.64 \\
      & $\mathcal{X}_\mathcal{P}$ & \phantom{0}49.14 & \phantom{0}96.20 & \phantom{0}96.00 & \phantom{0}92.14 \\
      \midrule
      & $\mathcal{X}_\mathcal{U}$ & \phantom{000}.00 & \phantom{000}.10 & \phantom{000}.06 & \phantom{000}.05 \\
      \textbf{Coverage}  & $\mathcal{X}_\mathcal{N}$ & \phantom{000}.00 & \phantom{000}.10 & \phantom{000}.06 & \phantom{000}.05 \\
      & $\mathcal{X}_\mathcal{P}$ & \phantom{00}7.33 & \phantom{0}60.79 & \phantom{0}63.88 & \phantom{0}78.48 \\
      \midrule
      & $\mathcal{X}_\mathcal{U}$ &           100.00 &           100.10 &           100.00 &           100.00 \\
      \textbf{Precision} & $\mathcal{X}_\mathcal{N}$ &           100.00 &           100.10 &           100.00 &           100.00 \\
      & $\mathcal{X}_\mathcal{P}$ & \phantom{0}80.83 & \phantom{0}97.76 & \phantom{0}99.29 & \phantom{0}99.10 \\
      \midrule
      & $\mathcal{X}_\mathcal{U}$ & \phantom{000}.00 & \phantom{000}.09 & \phantom{000}.06 & \phantom{000}.04 \\
      \textbf{Recall}    & $\mathcal{X}_\mathcal{N}$ & \phantom{000}.00 & \phantom{000}.10 & \phantom{000}.06 & \phantom{000}.05 \\
      & $\mathcal{X}_\mathcal{P}$ & \phantom{0}66.92 & \phantom{0}96.75 & \phantom{0}96.86 & \phantom{0}95.07 \\
      \bottomrule
    \end{tabular}
  \end{threeparttable}
\end{table}

\begin{foldable}
  As shown, $\mathcal{X}_\mathcal{U}$ and $\mathcal{X}_\mathcal{N}$ form a compact, separate cluster far from real images, indicating that their semantics are limited and dissimilar to natural objects.
  In contrast, $\mathcal{X}_\mathcal{P}$ overlaps with real image regions, suggesting that they share similar semantics.
  For quantitative analysis, we evaluate four distribution metrics (precision, recall, density, coverage) from~\cite{naeem2020reliable}.
  The statistics in Table~\ref{tab:abla-embeddings-coverage-objects} show that our image priors $\mathcal{X}_\mathcal{P}$ achieve significantly higher density, coverage, and recall at a small trade-off in precision, compared to noise sources $\mathcal{X}_\mathcal{U}$ and $\mathcal{X}_\mathcal{N}$.
  These results further confirm that $\mathcal{X}_\mathcal{P}$ has superior fidelity and diversity, and explain why $\mathcal{X}_\mathcal{P}$ excels in BBDFKD.
\end{foldable}

\subsubsection{Embeddings analysis --- Contrast} \label{sssec:04-experiments-ablation-embeddings-contrast}
\begin{foldable}
  To analyze the Contrast phase's impact on diversity, we evaluate 50K image prior embeddings in Fig.~\ref{fig:embeddings-cifar10}.
  The backbone embeds real images into ten clusters (equivalent CIFAR10 classes; gray).
  While pre-Contrast priors $\mathcal{X}_\mathcal{P}$ (green) already scatter around real images, the Contrast phase (orange) expands coverage into neighboring regions.
  This increased semantic diversity lets the student capture broader teacher knowledge.
  The distribution metrics in Table~\ref{tab:abla-embeddings-coverage-cifar10} confirm this with a significant increase in coverage and recall, with a trivial trade-off in density.
  Overall, these results show that through fulfilling the contrastive objective, the Contrast phase improves data diversity semantically and empirically.
\end{foldable}

\begin{figure}[ht] % fig:embeddings-cifar10
  \centering
  \includegraphics[trim={120 0 0 0},clip,width=0.90\columnwidth]{./figures/04-experiments-vizembed-2d-cifar10-resnet18.png}
  \caption{Feature embeddings of image priors $\mathcal{X}_\mathcal{P}$ before and after the Contrast phase.}
  \label{fig:embeddings-cifar10}
\end{figure}

\begin{table}[ht]
  \centering
  \begin{threeparttable}
    \caption{Distribution metrics (\%) during the Contrast phase.}
    \label{tab:abla-embeddings-coverage-cifar10}
    \begin{tabular}{l c c c c}
      \toprule
      & \textbf{Density} & \textbf{Coverage} & \textbf{Precision} & \textbf{Recall} \\
      \midrule
      Before Contrast & 70.46            & \phantom{0}6.13   & 91.69              & 82.18 \\
      After Contrast  & 66.93            &           14.66   & 91.68              & 94.00 \\
      \bottomrule
    \end{tabular}
  \end{threeparttable}
\end{table}

\subsubsection{Increasing synthetic dataset size} \label{sssec:04-experiments-ablation-nsamples}
\begin{figure}[H] % fig:abla-nsamples
  \centering
  \includegraphics[width=0.80\columnwidth]{./figures/experiments-abla-nsamples.pdf}
  \caption{Distillation accuracy (\%) by number of image priors $N$.}
  \label{fig:abla-nsamples}
\end{figure}

\begin{foldable}
  The relationship between synthetic data budget $N$ and KD performance mirrors the emergence of collective intelligence.
  While it is expected that KD performance improves with large $N$ (Fig.~\ref{fig:abla-nsamples}), we are interested in the optimal cut-off point, which indicates the \textit{minimum quorum} required for the student to replicate the teacher's expertise.
  This threshold scales directly with the structural complexity of the environment.
  For digits datasets (MNIST, USPS, SVHN), minimal data ($N$ = 10K) is required to reach good results, implying informational redundancy.
  In large-scale settings (CIFAR100, Tiny-ImageNet, Imagenette), the student continues to improve around $N$ = 100K, suggesting that the information density has not yet saturated and the collective wisdom of the teacher's priors remains beneficial.
\end{foldable}

\subsubsection{Cross-architecture KD} \label{sssec:04-experiments-ablation-crossarch}
\begin{table}[ht]
  \centering
  \caption{Distillation accuracy (\%) of cross-architecture KD.}
  \label{tab:abla-crossarch}
  \begin{threeparttable}
    \begin{tabular}{l l l l}
      \toprule
      & \textbf{FMNIST} & \textbf{SVHN} & \textbf{CIFAR10} \\
      \midrule
      Teacher  & LeNet5 & AlexNet & ResNet18 \\
      & 90.90 & 96.16 & 95.21 \\
      \midrule
      LeNet5   & \meanstdbf{83.98}{0.65} & \meanstd{64.44}{1.83} & \meanstd{32.01}{3.18} \\
      AlexNet  & \meanstd{78.98}{0.92} & \meanstdbf{95.94}{0.05} & \meanstd{61.30}{0.69} \\
      ResNet18 & \meanstd{81.69}{1.58} & \meanstd{95.91}{0.04} & \meanstdbf{83.41}{0.46} \\
      \bottomrule
    \end{tabular}
    \begin{tablenotes}[flushleft] \footnotesize
    \item {[\textbf{bold}] best results.}
    \end{tablenotes}
  \end{threeparttable}
\end{table}

\begin{foldable}
  An under-investigated setting in BBDFKD methods is when the teacher and student have different architectures.
  We evaluate KD with heterogeneous student architectures (LeNet5, AlexNet, ResNet18) on FMNIST, SVHN, and CIFAR10.
  Table~\ref{tab:abla-crossarch} shows that students perform best when matching the teacher's architecture (diagonal), with performance degrading as the capacity gap widens as expected.
  Thus, DIP-KD is practical in real-world scenarios with unknown teachers, provided the student has sufficient capacity.
\end{foldable}

\subsubsection{Students of compressed capacity} \label{sssec:04-experiments-ablation-compress}
\begin{figure}[ht] % fig:abla-compression
  \centering
  \includegraphics[width=0.80\columnwidth]{./figures/experiments-abla-compression.pdf}
  \caption{Distillation accuracy (\%) by compression ratio (CR).}
  \label{fig:abla-compression}
\end{figure}

\begin{foldable}
  Beyond same-size configurations, we evaluate DIP-KD under aggressive compression, a scenario unaddressed in prior methods~\cite{wang2021zero,zhang2022ideal,yuan2024data}.
  Specifically, we distill student networks with 25\%, 4\%, and 1\% compression ratio (CR) by reducing channel counts by 2$\times$, 5$\times$, 10$\times$, respectively.
  Results indicate strong empirical robustness: performance remains stable at 25\% CR, with only tolerable decline at 4\% CR.
  Significant degradation occurs only at 1\% CR, where the student's extreme sparsity prevents it from capturing the complex inter-class signals.
  These results demonstrate that DIP-KD is ideal for \textit{resource-constrained} mobile or edge environments.
\end{foldable}
